% EE 5290 final project report.
% Compile with:  pdflatex report.tex  (twice for cross-references)
\documentclass[10pt,letterpaper]{article}

\usepackage[margin=0.95in]{geometry}
\usepackage[utf8]{inputenc}
\usepackage{microtype}
\usepackage{graphicx}
\usepackage{booktabs}
\usepackage{array}
\usepackage{caption}
\captionsetup{font=small,labelfont=bf}
\usepackage{hyperref}
\hypersetup{colorlinks=true,linkcolor=blue!50!black,urlcolor=blue!60!black,citecolor=blue!50!black}
\usepackage{amsmath,amssymb}
\usepackage{xcolor}
\usepackage{enumitem}
\setlist[itemize]{leftmargin=1.2em,topsep=2pt,parsep=0pt,itemsep=2pt}
\setlist[enumerate]{leftmargin=1.5em,topsep=2pt,parsep=0pt,itemsep=2pt}
\usepackage{titlesec}
\titlespacing*{\section}{0pt}{1.2ex plus 0.5ex minus .2ex}{0.6ex plus .2ex}
\titlespacing*{\subsection}{0pt}{0.8ex plus 0.4ex minus .2ex}{0.4ex plus .2ex}
\titleformat*{\section}{\large\bfseries}
\titleformat*{\subsection}{\normalsize\bfseries}

\setlength{\parskip}{2pt}
\setlength{\parindent}{0pt}

\title{\vspace{-2em}Predictive Targeting and Segmentation of \\ Food-Insecure Households in the U.S. FoodAPS Survey}
\author{EE 5290 -- Final Project Report \\ Iowa State University, Spring 2026}
\date{May 14, 2026}

\begin{document}
\maketitle
\vspace{-2.5em}

\begin{abstract}
\small
Food-assistance organizations face increasing demand under constrained
funding and capacity.  Using a cleaned, household-level extract of the
USDA FoodAPS survey (4{,}826 households, 38 variables), we develop a
two-track data-analytics framework that helps such organizations
\emph{(i)} predict which households are food insecure and
\emph{(ii)} discover \emph{types} of food-insecure households whose profiles
suggest distinct outreach strategies.  We compare $L_2$-regularized
logistic regression, an RBF-kernel SVM, and a random-forest classifier
under three class-imbalance strategies, and we use PCA followed by
$K$-means clustering, both with and without geographic features, to
expose household personas.  On the held-out 20\% test set, logistic
regression achieves PR-AUC $0.568$ and ROC-AUC $0.783$, with a
top-decile precision of 67\% capturing 24\% of food-insecure
households.  Geography-free persona clustering identifies five
interpretable groups; 88\% of the model's top-decile flags fall into
two of them (\emph{non-working low-income singles} and \emph{large
low-income families with children}), suggesting that targeted outreach
can focus effort on programs designed for these distinct populations.
\end{abstract}

% --- 1. Introduction ----------------------------------------------------------
\section{Introduction}
Food security---reliable access to nutritious food---is a foundational
component of household well-being.  Across the Midwest in particular,
demand for assistance is rising while organisational capacity is not,
forcing food banks and outreach groups to make explicit choices about
\emph{which} households to prioritize.  This project asks two
related questions on the cleaned FoodAPS-derived dataset provided
for the case:

\begin{enumerate}
\item \textbf{Prediction.}  Given household characteristics, can we
      predict whether the household is food insecure with enough
      accuracy to be useful for triage?
\item \textbf{Segmentation.}  Among food-insecure households, what
      \emph{kinds} of households exist, and how should that affect
      outreach?
\end{enumerate}

We address both questions with library-first scikit-learn modelling
\cite{scikit-learn}, focusing depth on (i) a controlled comparison of
three classifiers under three class-imbalance strategies, (ii) a careful
choice of $K$ for $K$-means based on silhouette and stability, and
(iii) a synthesis section that joins the model's risk score to each
household's persona to produce concrete targeting recommendations.

\textbf{Contributions.}
(1) A reusable preprocessing pipeline for the FoodAPS extract that
recodes survey sentinel codes, encodes mixed nominal/ordinal features,
and adds log-transformed access variables;
(2) a head-to-head evaluation of logistic regression, RBF-kernel SVM
and random forest on imbalanced food-insecurity prediction with
operating-point analyses suited to a triage workflow;
(3) two complementary clusterings of the food-insecure subset---a
geography-driven view and a five-persona, geography-free view---and
(4) a synthesis showing how the persona mix changes across the model's
top-decile to top-half operating points.

\textbf{Non-causal framing.} The dataset is observational and a survey
sample.  All findings are reported as \emph{associations} between
features and the food-security outcome, never as causal effects.

% --- 2. Background ------------------------------------------------------------
\section{Background and Methods}
\label{sec:background}
We summarize the four learning methods used in the report; full
derivations are standard in the course textbooks
\cite{murphy2022,boyd2018,blum2020}.

\paragraph{Logistic regression.}
Given features $\mathbf{x}\in\mathbb{R}^d$ and a binary outcome
$y\in\{0,1\}$, the model
$\Pr(y=1\mid\mathbf{x})=\sigma(\mathbf{w}^\top\mathbf{x}+b)$ with
$\sigma(z)=1/(1+e^{-z})$ is fit by minimizing the negative
log-likelihood with $L_2$ regularization controlled by $C$.  The
coefficients $\mathbf{w}$ are interpretable on standardized
inputs: $w_j>0$ raises the predicted probability when feature $j$
increases.

\paragraph{Kernel SVM.}
The dual form of the support vector machine
\cite{murphy2022} replaces inner products $\mathbf{x}_i^\top\mathbf{x}_j$
with a kernel $k(\mathbf{x}_i,\mathbf{x}_j)$.  We use the radial basis
function (RBF) kernel
$k(\mathbf{x}_i,\mathbf{x}_j)=\exp(-\gamma\|\mathbf{x}_i-\mathbf{x}_j\|^2)$
with $\gamma$ controlling locality and $C$ the soft-margin penalty.
Probabilities are produced by Platt scaling on cross-validated
margins.

\paragraph{Random forest.}
Ensembles of decision trees fit on bootstrap samples with
feature-bagging give a strong nonlinear baseline that captures
interactions automatically and yields impurity-based feature
importances \cite{murphy2022}.

\paragraph{PCA and $K$-means.}
PCA computes the top $k$ singular vectors of the centered standardized
data matrix; the projection retains the directions of maximum variance
\cite{boyd2018}.  $K$-means partitions samples into $k$ clusters by
alternating assignment to the nearest centroid and recomputation of
centroids; we choose $k$ by inspecting the elbow of the within-cluster
sum-of-squares (inertia) jointly with the silhouette score (the
average ratio of between-cluster to within-cluster distance).

\paragraph{Class imbalance.}
With a positive-class rate near 28\%, accuracy is misleading: a
trivial all-negative classifier achieves 72\% accuracy.  We therefore
optimize and report \emph{average precision} (PR-AUC) and supplement
it with \emph{recall at fixed precision} thresholds.  We compare three
strategies: no adjustment, balanced class weights (cost-sensitive
loss), and probability-threshold tuning chosen on out-of-fold training
predictions.

% --- 3. Data ----------------------------------------------------------------
\section{Data and Preprocessing}
\label{sec:data}

The provided dataset
(\texttt{data/Case\_dataset.csv}) contains
4{,}826 households described by 38 columns covering food-security
outcomes, economic conditions, household composition, employment and
demographics, food-assistance program participation, transportation
and access indicators, and food-acquisition behavior.  The binary
outcome is \texttt{food\_insecure\_flag\_adult}; 27.85\% of the rows
are positive (Figure~\ref{fig:class}).  The four-level
\texttt{adltfscat} and the household-reported \texttt{foodsufficient}
are excluded from the feature matrix to avoid label leakage.

\subsection{Cleaning decisions}
Survey sentinel codes ($-996$, $-997$, $-998$) appear in
\texttt{foodpantry}, \texttt{anyvehicle}, \texttt{vehiclenum}, and
\texttt{caraccess}; we recode them to \texttt{NaN} before any
analysis.  Two columns are dropped entirely:
\texttt{caraccess} (92\% sentinel-coded) and
\texttt{fah\_storetype\_unique} (mean $\approx 0.002$, near-constant).
String-coded categorical heads (\texttt{head\_hispanic},
\texttt{head\_racecat}) carry the response code \texttt{R} for
``refused''; we recode \texttt{R} to \texttt{NaN}.  The survey weight
\texttt{hhwgt} is documented but not used in modelling
(see Section~\ref{sec:limits}).

\subsection{Feature engineering}
Distance variables (\texttt{dist\_sm}, \texttt{dist\_walmart},
\texttt{nearsnap\_dist} and similar) are right-skewed; we
\emph{add} a $\log(1+x)$ transform alongside the raw value, so that
linear models see a near-symmetric signal while tree models retain
the original units.  Four household-composition derivations are
added: children per adult, elderly share, child share, and employed-adult
share.  Nominal categoricals (\texttt{head\_sex}, \texttt{head\_hispanic},
\texttt{head\_racecat}, \texttt{head\_employment}, \texttt{region},
\texttt{poverty\_band}, \texttt{targetgroup}) are one-hot encoded with
\texttt{NaN} mapped to an explicit ``unknown'' level.  Numeric NaN are
median-imputed.  The resulting matrix has $63$ features after
encoding.

\subsection{Train/test split and EDA highlights}
We hold out a stratified 20\% of households as a final test set (966
rows) using \texttt{random\_state}=42.  All hyperparameter tuning uses
5-fold stratified CV on the training fold only; the test set is touched
exactly once at the end of Section~\ref{sec:track1}.

\begin{figure}[!htb]
\centering
\includegraphics[width=0.45\linewidth]{figures/01_class_balance.png}\hfill
\includegraphics[width=0.50\linewidth]{figures/02_marginal_rates.png}
\caption{(left) Class balance on \texttt{food\_insecure\_flag\_adult}: 27.9\%
positive rate.  (right) Marginal food-insecurity rates by selected
features.  Households at or below the poverty guideline have nearly a
50\% positive rate; the rate is also visibly elevated for SNAP-using
households and households without an employed adult.  Bars are
descriptive only and do not control for confounding.}
\label{fig:class}
\end{figure}

% --- 4. Track 1: prediction --------------------------------------------------
\section{Track 1: Predictive Modelling}
\label{sec:track1}

We compare three classifiers under three imbalance strategies for a
total of nine evaluated configurations.  For each base classifier we
run \texttt{GridSearchCV} over a small grid (LR: $C\in\{0.01,0.1,1,10\}$;
SVM-RBF: $C\in\{0.5,1,5\}$, $\gamma\in\{\text{scale},0.01,0.1\}$;
RF: \texttt{max\_depth}$\in\{\textsc{None},10,20\}$,
\texttt{min\_samples\_leaf}$\in\{1,5\}$ with $300$ trees) using 5-fold
stratified CV on the training fold and the average-precision (PR-AUC)
score function.

\subsection{Threshold-tuning protocol}
For each model, we use scikit-learn's \texttt{cross\_val\_predict} to
generate \emph{out-of-fold} probability estimates on the training rows.
We then choose the probability threshold that maximizes F1 on those
out-of-fold predictions.  This avoids using the test set to choose the
threshold, which would inflate the reported metrics.

\subsection{Test-set results}
Table~\ref{tab:9row} summarizes performance on the held-out 966 rows.
Three observations stand out:

\begin{enumerate}
\item \textbf{Logistic regression wins on every aggregate metric}:
      ROC-AUC $0.7831$ and PR-AUC $0.5681$.  RF and SVM-RBF reach
      PR-AUC $\approx 0.52$.  This is consistent with a problem whose
      signal is largely linear in the engineered features (a
      coincidence we exploit in Section~\ref{sec:track1-interp}).
\item \textbf{Class weights and threshold tuning are nearly equivalent
      and substantially better than the default 0.5 threshold} for the
      F1 / recall combination.  Both strategies trade roughly 17\% of
      precision for an extra 30--50 percentage points of recall.  The
      operationally relevant message is that the choice of threshold
      matters more than the choice of class-weight scheme.
\item \textbf{Recall at high precision is small for every method}: at
      a 70\% precision floor, the best model recovers only
      $\approx 26\%$ of the food-insecure population.  This sets a
      realistic expectation for the targeting recommendation in
      Section~\ref{sec:synth}.
\end{enumerate}

\begin{table}[!htb]
\centering
\small
\caption{Test-set performance for nine (model, imbalance-strategy) combinations.
``F1-tuned thr'' picks the threshold that maximizes F1 on training
out-of-fold predictions.  Bold rows are the per-model best F1.}
\label{tab:9row}
\begin{tabular}{lrrrrrrrr}
\toprule
Model (strategy) & thr & ROC-AUC & PR-AUC & Brier & Prec. & Rec. & F1 & R@P=0.50 \\
\midrule
LR -- no adjust              & 0.50 & 0.7831 & 0.5681 & 0.161 & 0.624 & 0.364 & 0.460 & 0.610 \\
LR -- balanced               & 0.50 & 0.7829 & 0.5681 & 0.195 & 0.463 & 0.762 & 0.576 & 0.598 \\
\textbf{LR -- F1-tuned thr}  & 0.289 & 0.7831 & 0.5681 & 0.161 & 0.466 & 0.755 & \textbf{0.576} & 0.610 \\
\midrule
SVM-RBF -- no adjust          & 0.50 & 0.7551 & 0.5290 & 0.171 & 0.549 & 0.208 & 0.302 & 0.565 \\
SVM-RBF -- balanced           & 0.50 & 0.7677 & 0.5269 & 0.166 & 0.564 & 0.346 & 0.429 & 0.651 \\
\textbf{SVM-RBF -- F1-tuned thr} & 0.236 & 0.7551 & 0.5290 & 0.171 & 0.425 & 0.747 & \textbf{0.542} & 0.565 \\
\midrule
RF -- no adjust               & 0.50 & 0.7667 & 0.5208 & 0.166 & 0.602 & 0.186 & 0.284 & 0.524 \\
RF -- balanced                & 0.50 & 0.7665 & 0.5101 & 0.179 & 0.485 & 0.665 & 0.561 & 0.658 \\
\textbf{RF -- F1-tuned thr}   & 0.301 & 0.7667 & 0.5208 & 0.166 & 0.453 & 0.792 & \textbf{0.577} & 0.524 \\
\bottomrule
\end{tabular}
\end{table}

\subsection{Visual diagnostics}
Figure~\ref{fig:pr-cal} (left) shows the precision/recall curves for
the three classifiers' raw probability outputs.  LR's curve dominates
SVM and RF in the high-recall regime, the regime that matters for a
food bank that wants to maximize coverage of the food-insecure
population subject to a precision floor.  Figure~\ref{fig:pr-cal}
(right) shows the calibration diagnostics: LR is well calibrated
across the range; the random forest is mildly over-confident at the
extremes; the SVM is well-behaved in the middle but slightly
under-confident at the low end (an artefact of Platt scaling on a
limited sample).

\begin{figure}[!htb]
\centering
\includegraphics[width=0.34\linewidth]{figures/05_pr_curves.png}\hfill
\includegraphics[width=0.62\linewidth]{figures/06_calibration.png}
\caption{(left) Test-set precision-recall curves.  The dashed
horizontal line is the random-baseline at the test-set positive
rate of 0.279.  (right) Quantile-binned reliability diagrams for
each classifier; closer to the diagonal is better.}
\label{fig:pr-cal}
\end{figure}

\subsection{Interpretation}
\label{sec:track1-interp}

\begin{figure}[!htb]
\centering
\includegraphics[width=\linewidth]{figures/08_feature_importance.png}
\caption{Top-20 features.  \emph{Left:} standardized logistic-regression
coefficients (positive raises the probability of food insecurity).
\emph{Right:} random-forest impurity-based importances.  Both methods
agree that food pantry use, poverty band, sampling target group, and
distances to retailers are the dominant signals; LR additionally
identifies the protective effect of vehicles, more elderly members,
and education.}
\label{fig:fi}
\end{figure}

The two interpretation views in Figure~\ref{fig:fi} are remarkably
consistent.  Food-pantry use, low poverty band, and the FoodAPS
sampling target group dominate both rankings; the random forest
additionally elevates the cluster of distance variables, reflecting
that geographical access is the largest source of nonlinear variation
even after controlling for income.  Logistic regression---which can
report \emph{direction}---shows expected protective effects of more
vehicles, higher education, and elderly household members (the latter
likely reflects greater enrollment in income-stabilising programs such
as Social Security and SNAP).  We use \texttt{foodpantry} as a
predictor because the data dictionary lists it under ``program
participation''; we acknowledge it is a behavioural correlate of
insecurity rather than an independent causal driver, and note that it
contributes meaningfully but not overwhelmingly to the model
(coefficient $\approx 0.85$, importance rank $15$).

% --- 5. Track 2: segmentation -----------------------------------------------
\section{Track 2: Segmentation of Food-Insecure Households}
\label{sec:track2}

The 1{,}344 food-insecure households are projected onto principal
components of the standardized 63-feature matrix; 24 components
explain 80\% of the variance.  We then sweep $K$-means for $k\in\{2,\ldots,8\}$
with 20 random restarts per $k$.

\subsection{Geographic clustering: a strong rural/urban axis}
The silhouette criterion identifies $k=2$ ($\text{silhouette}=0.221$) as
the clear winner; $k=3,\ldots,8$ all hover near $0.10$.  The two clusters
correspond cleanly to rural ($n=222$, mean
\texttt{dist\_walmart}\,$\approx\!11$ miles) and non-rural
($n=1{,}122$, mean \texttt{dist\_walmart}\,$\approx\!3$ miles)
households.  Stability across eight random initializations is high
(mean ARI $> 0.9$).  This is a real but uninformative finding for
outreach planning: it tells the food bank that physical access
dominates within-population variation but not what programmes to design
within each access tier.

\subsection{Geography-free persona clustering}
\label{sec:personas}
We therefore re-cluster after dropping the seven geographic features
(distances, rural, nonmetro, region one-hots).  The remaining 50
features capture household composition, employment, demographics,
poverty, program use, and food-acquisition behaviour.  Silhouette is
maximized at $k=2$ as before but more usable values appear at
$k\in\{3,4,5\}$; we cap at $k=5$ for narrative tractability and report
that solution.

\begin{figure}[!htb]
\centering
\includegraphics[width=\linewidth]{figures/14_persona_heatmap.png}
\caption{Geography-free persona clustering: column-wise
$z$-scores of selected features within each persona, computed
relative to the food-insecure subset's mean.  Red cells are above
the food-insecure mean; blue cells are below.}
\label{fig:personas}
\end{figure}

The five personas (Figure~\ref{fig:personas}) admit clean narrative
names:

\begin{itemize}
\item \textbf{P0 -- Working families just above poverty} ($n=266$, 19.8\% of the food-insecure subset).
       Nearly all have an employed adult; poverty ratio 1.43;
       moderate household size; pantry use only 9\%.
\item \textbf{P1 -- Non-working low-income singles, high pantry use} ($n=391$, 29.1\%).
       Small households (mean size 2.0); only 16\% have an employed
       adult; poverty ratio 0.85; pantry use 28\%---triple the food-insecure
       average.  This is the cluster most clearly served by direct food
       distribution.
\item \textbf{P2 -- Large low-income families with children} ($n=347$, 25.8\%).
       Mean household size 5.0 with 2.6 children; poverty ratio 0.85;
       partial employment; lower head education.  School-meal eligibility
       is highest in this cluster.
\item \textbf{P3 -- Working educated households near the boundary} ($n=214$, 15.9\%).
       Highest poverty ratio (3.05, well \emph{above} the poverty
       guideline) and highest education; near-universal employment.
       These are households whose food insecurity is unlikely to be
       resolved by income-eligibility programmes.
\item \textbf{P4 -- Elderly low-income households, high SNAP enrollment} ($n=126$, 9.4\%).
       Mean head age 70; lowest employment; highest SNAP enrolment
       (42\%).  Their inclusion in the food-insecure subset suggests
       SNAP is necessary but not always sufficient.
\end{itemize}

% --- 6. Synthesis -----------------------------------------------------------
\section{Synthesis: Targeting Recommendations}
\label{sec:synth}

We assign every household in the test set to a persona by passing it
through the same preprocessing $\to$ behaviour-only PCA $\to$
$K$-means pipeline used in Section~\ref{sec:personas}, and rank
households by the LR model's predicted probability.  This lets us
look at \emph{who} the model would actually be flagging at each
operating point.

\begin{figure}[!htb]
\centering
\includegraphics[width=0.55\linewidth]{figures/15_persona_mix_at_cutoffs.png}\hfill
\includegraphics[width=0.42\linewidth]{figures/16_precision_recall_vs_cutoff.png}
\caption{(left) Persona composition of the model's top-$X\%$
predicted-risk households on the test set; the legend names match
Section~\ref{sec:personas}.  (right) Precision and recall as a
function of the cutoff $X\%$.  Vertical dashed lines mark the 10\%,
20\%, and 30\% operating points discussed in the text.}
\label{fig:synth}
\end{figure}

\paragraph{Top-decile composition.}
At the top 10\% predicted-risk cutoff, precision is 68\% (97 of the
966 test households are flagged; 66 are truly food-insecure, capturing
24\% of the test set's food-insecure population).  Of those 97
flagged households, \textbf{45\% are P1 (non-working low-income singles)
and 43\% are P2 (large low-income families with children)}.  The two
clusters whose risk the model rates lowest---P3 (working educated
near boundary) and P4 (elderly on SNAP)---together account for under
2\% of the top-decile flags despite making up over half the population.

\paragraph{Operational implication.}
A food bank that adopts top-decile triage of the model can plan
outreach as if the workload would split evenly between two distinct
programme tracks: (a) supplemental food distribution for low-income
adult singles (potentially via existing pantry partnerships, since
this group already over-uses pantries); and (b) family-oriented
programmes for large low-income households with children, including
school-meal coordination and bulk-pack distributions.  The remaining
12\% of the top decile sits in P0 (working families just above
poverty) and warrants light-touch SNAP-enrolment outreach.

\paragraph{What the model misses.}
Persona P3 (working educated households just above the poverty line)
is 36.9\% of the population but only 15.9\% of the food-insecure
subset; the model rates them with mean probability 0.12 and they
account for $<\!0.5\%$ of the top-decile flags.  These households
may experience episodic insecurity that the survey's snapshot does
not capture well, and they are unlikely to be reached by
income-eligibility programmes.  This is a population that warrants
qualitative follow-up rather than algorithmic targeting.

\paragraph{Equity inspection.}
Persona P4 (elderly on SNAP) is under-represented in the top decile
relative to its share of food-insecure households (1.0\% versus
9.4\%).  Two interpretations are consistent with the data: the model
correctly de-prioritizes households whose insecurity is partly
already addressed by SNAP and Social Security; or the model
under-weights this group because the snapshot fails to capture
elderly-specific access barriers (mobility, dietary restrictions).
The recommendation is to monitor this group via existing SNAP
partnerships rather than treat the model's low scores as a green
light to disengage.

% --- 7. Limitations ---------------------------------------------------------
\section{Limitations and Future Work}
\label{sec:limits}

\textbf{Survey weights ignored.} The provided dataset is a probability
sample with weights in \texttt{hhwgt}.  We did not use them in fitting
or in evaluation, on the basis of the data dictionary's note that
weights are optional.  A weighted retrain---using \texttt{sample\_weight}
in the LR loss and a weighted scoring function---would produce
different coefficient magnitudes and may shift cluster sizes; we
recommend this as the first follow-up.

\textbf{Observational data.} Every coefficient and importance is an
\emph{association}, not a causal effect.  In particular, the
strong positive coefficient on \texttt{foodpantry} is best read as
``households that already use pantries are also reporting food
insecurity'', not ``pantry use causes food insecurity''.  The same
caveat applies to \texttt{snap\_any}: its protective sign here cannot
be read as a causal estimate of SNAP's effectiveness.

\textbf{Imbalance methods left unexplored.} We compared no adjustment,
balanced class weights, and threshold tuning, and found them roughly
equivalent in F1 with no method dominating across precision floors.
Future work should add SMOTE or its variants for synthetic minority
oversampling and cost-sensitive losses with explicit cost matrices, and
should evaluate them on a Bayesian decision-theoretic objective that
weights false-negative cost (missing a food-insecure household)
against false-positive cost (wasted outreach effort).

\textbf{Persona stability across outcome thresholds.} Our personas are
fit on households flagged on the binary outcome.  Refitting on the
4-level \texttt{adltfscat} (low vs.\ very low food security) and on
near-poverty households who narrowly avoided the binary flag would
test whether the personas track \emph{insecurity} or \emph{poverty}.

\textbf{External data.} The case allows supplementing with USDA
county-level access data and ACS demographics.  Adding even a single
county-level summary (e.g., median income, supermarket density) would
plausibly tighten the model and let us sanity-check distance
features against published food-desert maps.

\paragraph{Reproducibility.}  All results are reproduced by
\texttt{notebooks/01..04.ipynb} with \texttt{random\_state=42} and
\texttt{requirements.txt} pinned to scikit-learn 1.8.  See
\texttt{README.md} for setup.

% --- 8. Conclusion ---------------------------------------------------------
\section{Conclusion}
A library-first analysis of the FoodAPS extract is sufficient to
build both a useful predictive model (PR-AUC 0.57, top-decile
precision 67\%) and an interpretable five-persona segmentation of the
food-insecure population.  The synthesis of the two---reading off
\emph{who} the model flags at the top decile---turns an abstract
ranking into concrete programme recommendations: roughly half the
top decile is non-working low-income singles, half is large
low-income families with children, and the remaining 12\% is working
families on the poverty boundary.  This kind of analysis is exactly
the type of decision support that capacity-constrained food banks
need; the methodology generalizes to other household-survey datasets
with mixed nominal/ordinal features and an imbalanced binary
outcome.

% --- References ------------------------------------------------------------
\begin{thebibliography}{9}
\small
\bibitem{scikit-learn} F.~Pedregosa et al.  ``Scikit-learn: Machine Learning
in Python'', \emph{JMLR}, vol.~12, pp.~2825--2830, 2011.
\bibitem{murphy2022} K.~P.~Murphy.  \emph{Probabilistic Machine Learning:
An Introduction}.  MIT Press, 2022.
\bibitem{boyd2018} S.~Boyd and L.~Vandenberghe.  \emph{Introduction to
Applied Linear Algebra}.  Cambridge University Press, 2018.
\bibitem{blum2020} A.~Blum, J.~Hopcroft, and R.~Kannan.  \emph{Foundations
of Data Science}.  Cambridge University Press, 2020.
\bibitem{usda-foodaps} USDA Economic Research Service.  ``National
Household Food Acquisition and Purchase Survey (FoodAPS)''.
Accessed May 2026.
\end{thebibliography}

\end{document}
